\section{Conclusion}
\label{sec:conclusion}
ContractKV turns a structurally reusable prompt prefix into a governed multi-tenant KV service object. CALC exposes a legal physical prefix, scoped behavioral evidence determines its admissible use, PSTP controls first shared visibility, and a versioned Service-Template Certificate lets authorized consumers share qualification. The contribution lies in making these decisions one executable service lifecycle.

The measurements connect that lifecycle to useful service outcomes. Under task-preserving zero-slack query qualification, three models achieve 1.070--1.619$\times$ complete-service throughput, with identical raw output IDs on all 1,296 admitted future pairs. Fixed-layout template sharing reduces verifier calls from 125 to five and unhidden exposure from 25.05 to 3.24 ms/request; an evaluated open-arrival deployment point gains 0.658 SLO-compliant requests/s. The four-configuration complete-cost results in Table~\ref{tab:deployment-width} and the all-route physical authorization controls identify the costs and enforcement boundary behind those gains.

Three lessons follow: compilation creates reuse opportunity; evidence and transactional publication make it an admissible shared resource; and template-level qualification sharing, together with locality and residency, determines its service economics. Treating reusable KV as an evidence-bound service object provides a concrete basis for deploying transformed prompts as shared infrastructure.
